% =========================
% report/dataset.tex
% =========================

\subsection{توصیف داده}
دادهٔ آزمایشی شامل زوج‌های $(\mathcal{I}_T, \mathcal{I}_Q)$ است که در آن $\mathcal{I}_T$ تصویر الگو/نقشه و $\mathcal{I}_Q$ تصویر واقعی قطعه است.

\subsection{ملاحظات جمع‌آوری داده}
برای نتیجه‌گیری علمی معتبر، داده باید تغییرپذیری‌های زیر را پوشش دهد:
\begin{itemize}
  \item تغییر دید (Tilt/Rotation/Perspective)،
  \item تغییر نور (سایه، بازتاب، نور داخل/خارج)،
  \item شلوغی پس‌زمینه و عناصر مزاحم،
  \item انسداد جزئی (Occlusion) و قطع شدن قاب (Truncation).
\end{itemize}

\subsection{تقسیم‌بندی تنظیم/آزمون}
اگرچه مدل یادگیری وجود ندارد، اما تنظیم پارامترها می‌تواند باعث «بیش‌برازش تنظیمات» شود. بنابراین توصیه می‌شود:
\begin{itemize}
  \item یک زیرمجموعه برای تنظیم پارامترها (tuning),
  \item یک زیرمجموعهٔ مجزا برای گزارش نهایی (test)
\end{itemize}
در نظر گرفته شود.

\subsection{زمین‌حقیقت برای معیارهای کمی}
برای ارزیابی عددی دقیق، لازم است زمین‌حقیقت تعریف شود:
\begin{itemize}
  \item \textbf{نقاط متناظر} (Landmarks): چند نقطهٔ مشخص روی الگو و عکس به‌صورت دستی نشانه‌گذاری شود.
  \item \textbf{ماسک‌ها} (اختیاری): اگر هدف بررسی همپوشانی ناحیه‌ای باشد، ماسک دودویی قطعه می‌تواند IoU را ممکن کند.
\end{itemize}
در نبود زمین‌حقیقت، باید صراحتاً بیان شود که معیارهای جایگزین (مثل نسبت درون‌نهشتی‌ها) «نمایندهٔ مستقیم خطای واقعی» نیستند و صرفاً شاخص‌های تقریبی‌اند.
